\documentclass[11pt]{article}
\usepackage[a4paper,margin=25mm]{geometry}
\usepackage{amsmath,amssymb,amsthm,booktabs,array}
\usepackage[T1]{fontenc}
\usepackage{lmodern}
\usepackage[colorlinks=true,linkcolor=blue,urlcolor=blue,citecolor=blue]{hyperref}
\newtheorem{theorem}{Theorem}
\newtheorem{lemma}[theorem]{Lemma}
\newtheorem{corollary}[theorem]{Corollary}
\newcommand{\NN}{\mathbb{N}}
\title{The OEIS maximal-bishop growth conjecture:\\an extension of Arshon's enumeration}
\author{Haobo Ma\\{\small ChronoAI Pte Ltd, Singapore}
\and Wenlin Zhang\\{\small National University of Singapore}\\{\small ChronoAI Pte Ltd, Singapore}}
\date{November 17, 2026}
\begin{document}
\maketitle
\begin{abstract}
The formula at line 12 of the pinned export of OEIS entry A146304,
revision 35 (February 16, 2025),
conjectures that the number $a(n)$ of inclusion-maximal nonattacking bishop
placements on the fixed-coordinate $n\times n$ board is $O(n^{n-1})$;
we disprove this by extending Arshon's enumeration
(1936, pp.~25--29) \cite{arshon} to distinguishable maximal completions.
We construct independent placements whose distinct
maximal extensions retain their prescribed occupied diagonals.
For $n\geq3$ and $1\leq t\leq n-2$ we prove
$t^{2(n-1-t)}\leq a(n)\leq B_{n+2}^{\,2}$, where $B_j$ is the $j$th
Bell number. It follows that $a(n)^{1/(2n)}\sim n/(e\log n)$.
The same estimate holds for all nonattacking subsets; each color has the
corresponding one-color estimate. We also obtain growth estimates for
the largest coefficients of the associated counting polynomials and show
that a proportion tending to one of the maximal placements has
$2n-o(n)$ bishops.
\end{abstract}

\section{The counting problem}
Let $\mathcal{B}_n=\{0,\ldots,n-1\}^2$ for an integer $n\geq1$.
For a square $p=(x,y)$ write $d(p)=x-y$ and $s(p)=x+y$.
Distinct squares attack if their $d$-values or their $s$-values agree.
A placement is a subset of $\mathcal{B}_n$ containing no attacking pair.
It is \emph{maximal} if no further square can be added; this is different
from having maximum cardinality. Let $a(n)$ count all maximal placements,
with no identification under rotation or reflection.

The entry A146304 was submitted by Paolo Bonzini; its conjectured growth
formula has no separate signature. In the pinned export\footnote{\url{https://raw.githubusercontent.com/oeis/oeisdata/847798eafedfc5f486a00619e2926cfdb6bf73f9/seq/A146/A146304.seq}}
the conjecture is at line 12 and the counting convention at lines 5--6
\cite{oeis}. Andrew Howroyd's algorithm counts these placements through
the two board colors and gives a recurrence for their counting polynomials
\cite{howroyd}. We use the same color decomposition, but count a family of
partial placements whose completions are distinguishable. This gives direct
asymptotic bounds without evaluating the recurrence. A staircase-board
comparison supplies the logarithmic correction to the growth law.
Arshon's row-product rule already counts placements on nested rows
\cite[p.~25, equations (1)--(2)]{arshon}; we transfer such counts to
maximal extensions by recovering prescribed occupied diagonals.
Kaplansky and Riordan develop the classical rook framework
\cite{kaplansky}; Robinson studies prescribed-size dominating placements
\cite{robinson}. Arshon's cardinality-refined bishop enumeration is
reproved and simplified by Santos \cite{santos}. These earlier exact
enumerations underlie the related all-subset growth estimates; the
maximal-completion argument and the quantitative bounds are given below.

Let $A(n)$ count all nonattacking subsets, including the empty subset.
There is also a direct connection between the older enumerations and
the coarser normalized-logarithm limit. Every nonattacking subset
extends to a maximal one of size at most $2n-2$ for $n\geq2$.
Assigning each subset to one extension therefore gives
$A(n)\leq2^{2n-2}a(n)$: any fixed extension has at most $2^{2n-2}$
subsets. Thus an estimate $\log A(n)/(n\log n)\to2$ obtained from
the classical formulas also implies the same limit for $a(n)$.
The explicit prescribed-diagonal lower bound below gives the sharper
root equivalent $n/(e\log n)$, including the linear term in its logarithm.

Let $B_j$ count the set partitions of a $j$-element set, with $B_0=1$.
All logarithms below are natural.

\begin{theorem}\label{thm:main}
For every integer $n\geq3$ and every integer $t$ with $1\leq t\leq n-2$,
\begin{equation}\label{eq:bounds}
 t^{2(n-1-t)}\ \leq\ a(n)\ \leq\ A(n)\ \leq\
 \min\{(n+1)(n!)^2,B_{n+2}^{\,2}\}.
\end{equation}
In particular,
\begin{equation}\label{eq:limit}
 \begin{aligned}
 \log a(n)&=\log A(n)+o(n),\\
 \log A(n)&=2n(\log n-\log\log n-1)
   +O\!\left(\frac{n\log\log n}{\log n}\right).
 \end{aligned}
\end{equation}
\end{theorem}

In particular,
\[
 a(n)^{1/(2n)}\sim A(n)^{1/(2n)}\sim\frac{n}{e\log n},\qquad
 (a(n)/A(n))^{1/n}\to1.
\]
Also $\log a(n)/(n\log n)\to2$.
These statements do not assert $a(n)/A(n)\to1$.

\section{Prescribed diagonals and maximal extensions}
Color a square by $d(p)\bmod2$.
Since $s(p)\equiv d(p)\pmod2$, squares of different colors never attack.
For a fixed color, order the diagonals $d(p)=\delta$ by increasing length,
breaking ties arbitrarily, and denote their lengths by $L_0,L_1,\ldots$.

\begin{lemma}\label{lem:length}
Each color has at least $n-1$ difference diagonals, and
$h\leq L_h\leq h+2$
whenever its $h$th diagonal is defined, with indexing starting at zero.
\end{lemma}
\begin{proof}
The difference labels are $-(n-1),\ldots,n-1$, and parity partitions
these $2n-1$ consecutive integers into classes of sizes $n$ and $n-1$.
The diagonal of difference $\delta$ has length $n-|\delta|$.
Each length $\ell<n$ occurs twice, and length $n$ occurs once;
within one color, all the lengths have a fixed parity.
For $h\geq1$, among the integers $1,\ldots,h-1$ at most
$\lceil(h-1)/2\rceil$ have that parity. Thus at most
$2\lceil(h-1)/2\rceil\leq h$ diagonals have length less than $h$.
The diagonal with index $h$ is the $(h+1)$st one, so its length is at
least $h$. The assertion for $h=0$ is immediate.
The length lists have the form $1,1,3,3,\ldots$ or
$2,2,4,4,\ldots$, with a possible unpaired final length, so
$L_h\leq h+2$.
\end{proof}

\begin{lemma}\label{lem:completion}
Fix a collection of difference diagonals. For each nonattacking partial
placement consisting of exactly one square on each of these diagonals, choose
any maximal extension. Distinct partial placements have distinct extensions.
\end{lemma}
\begin{proof}
A maximal extension exists: repeatedly add any unattacked square until none
remains, which terminates because the board is finite.
An extension cannot contain a second square on a prescribed diagonal.
Restricting the extension to the union of the prescribed diagonals therefore
recovers the partial placement. If two extensions were equal, their
restrictions would be equal as well.
\end{proof}

\begin{lemma}\label{lem:bell}
The number of all nonattacking placements on either color is at most
$B_{n+2}$. As $N\to\infty$,
\[
 \log B_N\leq N(\log N-\log\log N-1)
       +O\!\left(\frac{N\log\log N}{\log N}\right).
\]
\end{lemma}
\begin{proof}
Represent a square by the pair consisting of its difference and sum
diagonals. A nonattacking placement becomes a rook placement with at
most one square in each row and column. In a fixed color, the sum labels
on the difference diagonal $\delta$ are
\[
 |\delta|,\ |\delta|+2,\ldots,2n-2-|\delta|.
\]
These sets are nested when the difference diagonals are ordered by
length. Thus the columns can be reordered so that row $h$ occupies the
first $L_h$ columns. If there are $r\leq n$ rows, the bound
$L_h\leq h+2$ embeds this board into the staircase of row lengths
$1,2,\ldots,r+1$, by leaving its first row unused.

For completeness, write $R(m,k)$ for the number of placements of $k$
rooks on the staircase of lengths $1,\ldots,m$.
Deleting its last row gives
\[
 R(m,k)=R(m-1,k)+(m-k+1)R(m-1,k-1).
\]
The second term chooses the last-row rook after the other $k-1$ rooks;
exactly $m-k+1$ columns remain. With $R(0,0)=1$, this is the recurrence
for $R(m,k)=S(m+1,m+1-k)$, where $S(u,v)$ is the number of partitions
of a $u$-element set into $v$ nonempty blocks. Therefore the total is
$B_{m+1}$. The embedding gives $B_{r+2}\leq B_{n+2}$.
This staircase identity is part of the classical rook framework
\cite{kaplansky}; its proof here fixes the indexing.

Decomposing a set partition into its blocks gives
$\sum_{N\geq0}B_Nz^N/N!=\exp(e^z-1)$.
For real $z>0$, positivity of the coefficients implies
\[
 \log B_N\leq\log(N!)+e^z-1-N\log z.
\]
Take $z=\log N-\log\log N$ for sufficiently large $N$.
Then $e^z=N/\log N$ and
$\log z=\log\log N+O(\log\log N/\log N)$.
The integral bounds on $\sum_{j=1}^N\log j$ give
$\log(N!)=N\log N-N+O(\log N)$, proving the estimate.
\end{proof}

\begin{proof}[Proof of Theorem~\ref{thm:main}]
Fix $n,t$ in the stated range. In each color retain the $n-1$ shortest
difference diagonals and omit the first $t$ of these. Select one square on
each remaining diagonal, in the inherited order. At step $j$, starting
from $j=0$, the current diagonal has index $t+j$ and hence at least $t+j$
squares by Lemma~\ref{lem:length}.

On a fixed difference diagonal, $s(x,y)=x+y$ is injective: knowing $s$ and
$d$ determines $x=(s+d)/2$ and $y=(s-d)/2$.
Each of the $j$ previously chosen squares therefore forbids at most one
square on the current diagonal. At least $(t+j)-j=t$ choices remain.
There are $n-1-t$ steps, giving at least $t^{n-1-t}$ partial placements
of that color. Choices in different colors combine freely, giving at least
$t^{2(n-1-t)}$ partial placements on the board.
Lemma~\ref{lem:completion} turns these into distinct maximal placements.

For the upper bound, every nonattacking placement, whether maximal or not,
is determined by recording either its unique square or its absence on each
difference diagonal. Ignoring the restrictions from the sum diagonals only
increases the number of possible records. A diagonal of length $\ell$
has $\ell+1$ possible records. Their product is
\[
 (n+1)\prod_{\ell=1}^{n-1}(\ell+1)^2=(n+1)(n!)^2.
\]
The two color restrictions of any nonattacking subset are independent
choices, so Lemma~\ref{lem:bell} gives $A(n)\leq B_{n+2}^{\,2}$.

Finally, for sufficiently large $n$ take $t=\lfloor n/\log n\rfloor$,
which lies between $1$ and $n-2$.
Put $L=\log n$ and $\ell=\log\log n$. Since
$t=n/L+O(1)$ and $\log t=L-\ell+O(L/n)$, the lower bound gives
\[
 \log a(n)\geq2(n-1-t)\log t
 =2n(L-\ell-1)+O(n\ell/L).
\]
Lemma~\ref{lem:bell}, with $N=n+2$, gives the matching upper bound
for $\log A(n)$. Hence both logarithms have the stated expansion;
their difference is $O(n\ell/L)=o(n)$.
\end{proof}

\begin{corollary}\label{cor:refutation}
For every integer $m\geq1$, $a(4m+1)\geq m^{6m}$.
For every integer $m\geq50$,
\begin{equation}\label{eq:ratio}
 \frac{a(4m+1)}{(4m+1)^{4m}}\geq4^m.
\end{equation}
In particular, $a(n)$ is not $O(n^{n-1})$.
\end{corollary}
\begin{proof}
Take $n=4m+1$ and $t=m$ in \eqref{eq:bounds}.
For $m\geq50$, $m^2\geq2500$ and hence
\[
 4(5m)^4=2500m^4\leq m^6.
\]
Using $4m+1\leq5m$ and raising this inequality to the $m$th power gives
\[
 4^m(4m+1)^{4m}\leq4^m(5m)^{4m}
   =\bigl(4(5m)^4\bigr)^m\leq m^{6m}\leq a(4m+1).
\]
To see that this disproves the full eventual bound, let $C>0$ and let
$N\geq1$ be any cutoff. Choose $m\geq\max(50,N)$ with $4^m>C$.
Then $n=4m+1\geq N$, and \eqref{eq:ratio} yields
$a(n)>Cn^{n-1}$.
\end{proof}

\section{Related counts and verification scope}
Call the color containing $(0,0)$ black. Write $b(n),w(n)$ for its
black and white maximal-placement counts, and $U_b(n),U_w(n)$ for all
placements in those colors. A color with no squares has one placement,
the empty set.

\begin{corollary}\label{cor:colors}
Howroyd's identity $a(n)=b(n)w(n)$ holds. For
$q(n)\in\{b(n),w(n),U_b(n),U_w(n)\}$,
\[
 \log q(n)=n(\log n-\log\log n-1)
       +O\!\left(\frac{n\log\log n}{\log n}\right),
 \qquad q(n)^{1/n}\sim\frac{n}{e\log n}.
\]
The same counts enumerate minimal vertex covers, respectively all vertex
covers, in the corresponding bishop graphs.
\end{corollary}
\begin{proof}
A placement is maximal exactly when both color restrictions are
maximal: an addable square attacks nothing of the other color. This
proves the product identity \cite{howroyd}. The one-color construction
gives $t^{n-1-t}\leq b(n),w(n)\leq U_b(n),U_w(n)\leq B_{n+2}$.
The preceding logarithmic estimates apply with the factor $2$ removed.
In any graph, complementation bijects independent sets with vertex
covers and maximal independent sets with minimal vertex covers.
The parenthetical ``minimal edge covers'' in A290613, name line~5
of its pinned export,\footnote{\url{https://raw.githubusercontent.com/oeis/oeisdata/847798eafedfc5f486a00619e2926cfdb6bf73f9/seq/A290/A290613.seq}}
must therefore read ``minimal vertex covers''.
For $n=2$ the white bishop graph is a single edge: it has two maximal
independent sets and two minimal vertex covers, but only one minimal
edge cover, namely that edge. This corrects the parenthetical, not the
sequence's maximal-independent-set counts.
\end{proof}

For $n\geq2$ each color supports at most $n-1$ bishops. The color with
$n-1$ difference diagonals has this bound directly. In the other color,
the two extreme difference diagonals are singletons with the same sum,
so not all its $n$ diagonals can be occupied. Thus the full bound is
$2n-2$. The one-square boundary is $a(1)=1$.

Let $p_c(n,k)$ and $r_c(n,k)$ count the maximal, respectively all,
$k$-bishop placements in color $c\in\{b,w\}$. Define
$P_c(n,z)=\sum_k p_c(n,k)z^k$ and
$R_c(n,z)=\sum_k r_c(n,k)z^k$.
Write $P(n,z)=P_b(n,z)P_w(n,z)$ and
$R(n,z)=R_b(n,z)R_w(n,z)$ for the full-board polynomials.
For $n\geq2$ the color polynomials have degree at most $n-1$.

\begin{corollary}\label{cor:coefficients}
The row sums and the largest coefficients of each of $P_c(n,z)$ and
$R_c(n,z)$ have $n$th root asymptotic to $n/(e\log n)$.
The largest coefficients of $P(n,z)$ and $R(n,z)$ have $2n$th root
asymptotic to $n/(e\log n)$.
Every size attaining a largest full-board coefficient is $2n-o(n)$;
the corresponding color sizes are $n-o(n)$.
For every fixed $0<\varepsilon<1$, the proportion of maximal placements
with fewer than $(2-\varepsilon)n$ bishops tends to zero.
The same is true among all nonattacking placements.
\end{corollary}
\begin{proof}
For a polynomial with $d+1$ nonnegative coefficients, its largest
coefficient lies between its value at $1$ divided by $d+1$ and its
value at $1$. Here $d\leq n-1$ for a color and $d\leq2n-2$ for the
full board, so the logarithmic difference is $O(\log n)$.
Theorem~\ref{thm:main} and Corollary~\ref{cor:colors} give the estimates.

A $k$-bishop placement chooses $k$ of the $2n-1$ difference diagonals
and one of at most $n$ squares on each. Consequently the number with
at most $\lfloor(2-\varepsilon)n\rfloor$ bishops is at most
$2^{2n-1}n^{\lfloor(2-\varepsilon)n\rfloor}$.
Dividing by either $a(n)$ or $A(n)$ gives a logarithm at most
$-\varepsilon n\log n+2n\log\log n+O(n)$, which tends to
$-\infty$. Before division, the bound is also $o(a(n)/n)$ and
$o(A(n)/n)$, so no such size can attain a largest coefficient.
For a color the corresponding numerator is at most
$2^n n^{\lfloor(1-\varepsilon)n\rfloor}$, giving the color assertion.
\end{proof}

\begin{corollary}\label{cor:fixedsize}
The $n$-bishop count A122749 is $o(a(n))$, and the unrestricted
$n$-bishop count A002465 is $o(A(n))$.
The number of maximum-size placements is also $o(a(n))$.
\end{corollary}
\begin{proof}
Each numerator is at most $2^{2n-1}n^n$ by the same diagonal record.
The logarithmic estimates for the denominators prove both limits.
The exact A122749 formulas are already known \cite{robinson};
these limits do not replace their fixed-size enumeration.
Robinson also gives $2^n$ maximum-size placements for $n\geq2$
\cite[Section 4, p.~213]{robinson}, and $2^n/a(n)\to0$ by
Theorem~\ref{thm:main}.
\end{proof}

\paragraph{A general form of the construction.}
In any finite bipartite graph, choose distinct left vertices
$u_0,\ldots,u_{k-1}$ with degrees $d_0,\ldots,d_{k-1}$ and suppose
$d_j\geq j+1$. There are at least
$\prod_{j=0}^{k-1}(d_j-j)$ maximal matchings. Indeed, sequentially
match $u_j$ to an unused neighbor; at most $j$ choices are forbidden.
Extend each matching maximally. Its edge at every $u_j$ is unchanged,
so restriction to these vertices recovers the seed. The bishop
construction is this principle on the diagonal-incidence graph.

\begin{center}
\small
\begin{tabular}{@{}p{0.25\linewidth}p{0.65\linewidth}@{}}
\toprule
Source item & Coverage \\
\midrule
A146304, formula line & Refuted-and-corrected:
Theorem~\ref{thm:main} and Corollary~\ref{cor:refutation}. \\
A290594; A290613 & Not posed; proved: color growth in Corollary~\ref{cor:colors}.
The product identity is reproved; prior proof Howroyd, 2017. \\
A290613, name line 5 & Refuted-and-corrected: ``minimal edge covers''
is replaced by minimal vertex covers; the $n=2$ graph gives counts
$1$ and $2$, respectively (Corollary~\ref{cor:colors}). \\
A288182; A288183 & Not posed; proved: growth of row sums and largest coefficients in
Corollary~\ref{cor:coefficients}. Other individual coefficients are not
evaluated. \\
A201862 & Not posed; proved: all-subset growth in Theorem~\ref{thm:main};
the empty set is included. \\
A274105; A274106; A216078; A216332 & Not posed; proved: growth of color row sums and largest
coefficients, with the last two entries evaluated at $n+1$.
The exact counts have prior formulas; see \cite{santos}. \\
A378590 & Not posed; proved: growth of full-board row sum and largest coefficient in
Theorem~\ref{thm:main} and Corollary~\ref{cor:coefficients}. \\
A122749; A002465 & Not posed; proved: negligible fixed-$n$-bishop fractions in
Corollary~\ref{cor:fixedsize}; prior fixed-size enumerations cited. \\
Robinson, Section 4 & Not covered: exact intermediate-size enumeration.
Corollary~\ref{cor:coefficients} gives peak growth and concentration,
not asserted there. His $2^n$ formula gives the maximum-size fraction in
Corollary~\ref{cor:fixedsize}. \\
A146303 & Not covered: maximal queen counts. The seed extension does
not enforce distinct rows and columns. \\
\bottomrule
\end{tabular}
\end{center}

\paragraph{Formal statements.}
In the pinned
\href{https://github.com/the-omega-institute/trureturing/blob/549945461571e5d29c740b71d3b02651c85044f5/D5/S3/Combinatorics/Graph/Domination/BishopMaximalPlacementGrowth.lean}{\texttt{BishopMaximalPlacementGrowth.lean}},
the declarations \texttt{count\_lower\_bound},
\texttt{ratio\_lower\_bound}, and \texttt{result} formalize,
respectively, $m^{6m}\leq a(4m+1)$, \eqref{eq:ratio}, and the failure
of the conjectured eventual Big-O bound. The declarations
\texttt{physical\_cap} and \texttt{source\_family\_card} give, respectively,
the bound $2n-2$ for each maximal set when $n\geq2$, and equality of the
cap-filtered count with the native maximal-set count.
The all-$n$ parameterized lower bound, Bell and record upper bounds,
logarithmic expansions, coefficient estimates, concentration, and
general bipartite-graph statement have written proofs above.

\paragraph{Finite checks.}
The accompanying \texttt{verify.py} uses exact integer arithmetic to enumerate
the maximal placements of each color, giving
\[
 a(1),\ldots,a(8)=1,4,10,64,660,7744,111888,1960000,
\]
in agreement with A146304. It also checks the color coefficient tables,
all-subset counts, staircase identity, diagonal nesting, seed recovery,
and the arithmetic in \eqref{eq:ratio}. The exact ranges and outputs are
in \texttt{verification.json}. The infinite conclusions follow from the
written proofs.

\begin{thebibliography}{9}
\bibitem{oeis}
P. Bonzini and the OEIS contributors,
\emph{A146304: Number of distinct ways to place bishops on an $n\times n$
chessboard so that no bishop is attacking another and it is not possible to
add another bishop}, The On-Line Encyclopedia of Integer Sequences,
submitted October 29, 2008, revision 35, February 16, 2025.
\url{https://oeis.org/A146304}.
\bibitem{howroyd}
A. Howroyd, \emph{Number of maximal independent sets in $n\times n$
bishops graph and other rook-like graphs}, algorithm note linked from
A146304; associated program dated June 5, 2017.
\url{https://oeis.org/A146304/a146304.txt}.
\bibitem{arshon}
S. E. Arshon, \emph{Solution of one combinatorial problem} (in Russian),
Matematicheskoe Prosveshchenie, series 1, issue 8 (1936), 24--29;
row-product rule p.~25, factorial expansion p.~26, bishop formulas
pp.~27--29.
\url{https://www.mathnet.ru/eng/mp694}.
\bibitem{kaplansky}
I. Kaplansky and J. Riordan, \emph{The problem of the rooks and its
applications}, Duke Mathematical Journal \textbf{13} (1946), 259--268,
especially Section 9.
\url{https://doi.org/10.1215/S0012-7094-46-01324-5}.
\bibitem{robinson}
R. W. Robinson, \emph{Counting arrangements of bishops},
Combinatorial Mathematics IV, Lecture Notes in Mathematics
\textbf{560}, Springer, 1976, 198--214.
\url{https://doi.org/10.1007/BFb0097382}.
\bibitem{santos}
E. G. Santos, \emph{Counting Nonattacking Chess Piece Placements:
Bishops and Anassas}, Journal of Integer Sequences \textbf{28}
(2025), Article 25.8.6, December 15, 2025, Theorems 3--4, pp.~3--6.
His color names are interchanged relative to the black-corner convention
here. Earlier version: arXiv:2411.16492v2, August 13, 2025,
Theorems 2.1--2.2.
\url{https://cs.uwaterloo.ca/journals/JIS/VOL28/Santos/santos3.pdf}.
\end{thebibliography}
\medskip
\noindent\small\copyright\ 2026 Haobo Ma and Wenlin Zhang.
Licensed under Creative Commons Attribution 4.0 International (CC BY 4.0).
\end{document}
